\documentclass[11pt]{article}

\usepackage[margin=1in]{geometry}
\usepackage[T1]{fontenc}
\usepackage[utf8]{inputenc}
\usepackage{booktabs}
\usepackage{longtable}
\usepackage{array}
\usepackage{xcolor}
\usepackage{hyperref}
\usepackage{listings}
\usepackage{amsmath}
\usepackage{siunitx}

\hypersetup{
  hidelinks,
  pdftitle={GICForge Manual},
  pdfauthor={Merlino Development Team}
}

\definecolor{MerlinoBlue}{HTML}{245C73}
\definecolor{MerlinoSlate}{HTML}{4A5561}
\definecolor{MerlinoGold}{HTML}{C88A2D}
\definecolor{MerlinoBg}{HTML}{F7F8F6}
\emergencystretch=2em

\lstdefinestyle{merlino}{
  basicstyle=\ttfamily\small,
  breaklines=true,
  columns=fullflexible,
  keepspaces=true,
  frame=single,
  rulecolor=\color{MerlinoSlate!45},
  backgroundcolor=\color{MerlinoBg},
  xleftmargin=0.5em,
  xrightmargin=0.5em
}

\newcommand{\program}{\texttt{GICForge}}
\newcommand{\merlino}{\texttt{Merlino}}
\newcommand{\xyzin}{\texttt{xyzin}}
\newcommand{\provin}{\texttt{provin}}
\newcommand{\provout}{\texttt{provout}}
\newcommand{\gauin}{\texttt{gauin}}
\newcommand{\gicsym}{\texttt{gicsym}}
\newcommand{\angstrom}{\text{\AA}}

\title{\textbf{GICForge Manual}\\
\large Non-redundant generalized internal coordinates and symmetry-adapted coordinate bases in Merlino}
\author{Merlino Development Notes}
\date{\today}

\begin{document}
\maketitle

\begin{abstract}
\program{} is the \merlino{} geometry backend that converts Cartesian molecular
structures into non-redundant generalized internal coordinates (GICs), optional
symmetry-adapted GICs, and symmetry-adapted Cartesian coordinates.  It is used
as a library service by MORPHEUS, SEfit and future geometry optimizers, and can
also be run as a standalone diagnostic tool.  This manual describes the input
contract, coordinate-generation policy, pruning and symmetrization rules, output
files, Python/Fortran identity checks and recommended regression tests.
\end{abstract}

\tableofcontents
\newpage

\section{Purpose}

\program{} has one responsibility: starting from a Cartesian molecular
structure, it constructs a chemically interpretable, rank-complete coordinate
basis.  The basis can be used directly as a Gaussian ReadGIC block, frozen as a
machine-readable \merlino{} schema, evaluated analytically on later Cartesian
geometries, or converted into symmetry-adapted Cartesian displacements.

\program{} deliberately does not own project management, SMILES handling,
electronic-structure parsing, DVR, semiexperimental least-squares logic or GUI
workflow orchestration.  Those tasks belong to the Python layer.  The geometry
backend owns topology perception from Cartesian coordinates, primitive internal
coordinates, ring handling, type-local redundancy removal, symmetry labels and
the final auditable coordinate definition.

\section{Theoretical Background}

Let \(x\) be the \(3N\)-dimensional Cartesian coordinate vector and let
\(q_p(x)\) be a primitive internal coordinate.  \program{} first builds a
redundant primitive set
\[
  \mathbf{q}(x)=\{q_1(x),q_2(x),\ldots,q_M(x)\},
\]
where \(M\) is usually larger than the vibrational rank.  A generalized
internal coordinate is a linear combination of primitives,
\[
  Q_k(x)=\sum_p U_{pk} q_p(x),
\]
with \(U\) stored in the frozen GIC definition.  The Wilson row associated with
the coordinate is
\[
  \mathbf{B}_k = \frac{\partial Q_k}{\partial x}
              = \sum_p U_{pk}\frac{\partial q_p}{\partial x}.
\]
The final coordinate set is valid only if the B matrix has rank \(3N-6\) for a
nonlinear molecule or \(3N-5\) for a linear molecule.

The central algorithmic choice is that redundancy removal is performed by
coordinate family, not by a global decomposition over all candidates.  In a
global SVD or global rank-revealing decomposition, the numerical ordering of
rows can make chemically unrelated coordinates compete with one another.  A
stretch row could remove a torsional row, or an out-of-plane row could remove a
ring bend, even though the two coordinates describe different physical
motions.  \program{} therefore uses the already accepted rows as a rank basis
and admits additional rows only within a fixed family order.  Inside each
family, a modified Gram--Schmidt rank test keeps the rows that add linearly
independent B-matrix information.  This preserves interpretability while still
enforcing the exact vibrational rank.

Symmetry adaptation is a second projection problem.  Once the final
post-pruning basis is known, equivalent coordinate rows are projected with the
molecular symmetry operations and labelled by irreducible representation.  The
projection is constrained to homogeneous coordinate blocks, so the symmetry
operation cannot manufacture a coordinate by mixing unrelated physical
families.  The theoretical target counts are those of the vibrational
representation,
\[
  \Gamma_\mathrm{vib} = \Gamma_{3N}-\Gamma_\mathrm{trans}-\Gamma_\mathrm{rot}.
\]
If the projected coordinate counts cannot reproduce this representation, the
coordinate generator fails rather than writing an ambiguous basis.

\section{Input Contract}

The molecular geometry input is the XYZ-format file \xyzin{} in the working
directory.  Atomic symbols and atomic numbers must be equivalent and must cover
the full periodic table through the shared \merlino{} periodic-table utilities.
The \provin{} file contains only keywords, title, charge and multiplicity; it
must not contain Cartesian coordinates, Z-matrices or FCHK-derived geometry.

\begin{lstlisting}[style=merlino]
# GNIC BMAT GICSYM

Pyrene GICForge example

0 1
\end{lstlisting}

Unsupported legacy geometry paths are intentionally disabled: Z-matrix input,
Cartesian coordinates embedded in \provin{}, FCHK geometry input, SMILES input,
and old FITPOT/VCI/DVR/MSR/isotope/rate driver keywords.  Cartesian input is the
only production path.

\section{Build and Execution}

The active Fortran backend is built with:

\begin{lstlisting}[style=merlino]
cd fortran/gicforge
./compile_MAC
\end{lstlisting}

\noindent
The build writes \texttt{fortran/gicforge/build/gicforge} and installs the
runtime executable as \texttt{bin/gicforge.x}.  The compatibility alias
\texttt{bin/prova.x} is kept only for older launchers.

The recommended command-line entry point is the shared \merlino{} CLI:

\begin{lstlisting}[style=merlino]
python -m merlino gic --workdir work/gicforge
python -m merlino gic-define geometry.xyz --workdir work/gicdefine
python -m merlino gic-bmatrix merlino.gic.definition.json geometry.xyz
python -m merlino gic-contract --geometry geometry.xyz --workdir work/contract
\end{lstlisting}

\noindent
\texttt{gic-define} creates a frozen schema; \texttt{gic-bmatrix} evaluates that
schema without rerunning topology perception; \texttt{gic-contract} verifies the
Fortran/Python identity contract.

\section{Coordinate Construction Algorithm}

The coordinate generator is a deterministic pipeline.  Each step consumes only
information created by earlier steps, and the final schema records enough
metadata to verify that the same model is reused later.

\subsection{Topology perception}

\program{} starts from interatomic distances and atomic covalent radii.  The
result is a covalent graph with atoms as vertices and bonds as edges.  The graph
is checked before coordinate construction: symbols and atomic numbers must be
consistent, atom indexes must be unique, and obvious nonbonded H--H contacts or
other impossible short contacts are treated as input/topology errors.  The graph
is then frozen for the run.  Later refinement or optimization modules must not
allow a geometry step to change this graph silently.

Ring perception is performed on the covalent graph.  The ring basis is
canonicalized before any ring coordinate is built, so the same molecule gives
the same ring order in Fortran and Python.  Fused systems are represented by
their individual rings plus shared bonds; the shared bonds are the places where
butterfly coordinates can be introduced.

\subsection{Primitive coordinate generation}

The first coordinate layer is primitive and usually redundant.  The generator
adds primitives in chemically ordered families:

\begin{enumerate}
  \item stretches for every covalent edge;
  \item ordinary valence angles around each bonded center;
  \item linear-bend candidates only for nearly linear centers, with a maximum
        number per center;
  \item exocyclic torsions from the default one-dihedral rule;
  \item endocyclic valence-angle primitives for each ring;
  \item endocyclic torsional primitives for each ring;
  \item butterfly primitives around bonds shared by two rings;
  \item out-of-plane primitives, except when all four involved atoms are cyclic;
  \item optional improper-dihedral primitives in the Gaussian-compatibility mode.
\end{enumerate}

This order is not only cosmetic.  It defines the chemical priority used by the
rank reduction.  Stretches define the bonded frame first; bends define local
coordination; torsions and ring coordinates define conformational and
puckering directions; out-of-plane coordinates close the remaining local
non-planarity directions.

\subsection{Local combinations before pruning}

Several local coordination patterns are better represented by combinations than
by a single primitive.  For example, equivalent valence angles around a highly
symmetric center may be converted to local sums and differences, and ring
torsions are converted to deformation coordinates.  These combinations are
still type-local: angles are combined with angles, torsions with torsions, and
out-of-plane coordinates with out-of-plane coordinates.  This rule is essential
because it prevents the construction stage from creating mathematically valid
but chemically opaque coordinates.

For ring systems, endocyclic torsions are not meant to appear as an arbitrary
list of raw dihedrals in the final model.  They are the primitive source from
which \program{} builds the active \(q\)- and \(\phi\)-like ring deformation
coordinates.  The report therefore names the endocyclic torsional content in
the GIC definition section, but the final Gaussian block should contain the
selected \texttt{QPck} and \texttt{PhiP} coordinates after pruning and symmetry
assignment.

\subsection{Rank reduction and final basis}

After primitive generation and local combinations, \program{} evaluates Wilson
B rows for the candidate coordinates at the current Cartesian geometry.  It then
keeps rows in the fixed family order and rejects rows that do not increase the
rank.  The rank test is numerical, but the candidate competition is chemical:
an angle can remove another redundant angle, but not a stretch; a ring bend can
be removed by ring-bend redundancy, but not by an unrelated exocyclic torsion.

If the ordinary candidate set is too small before pruning, the backend expands
within Fortran to the full primitive candidate set and repeats the same
type-local pruning.  The number of pre-pruning attempts is capped by the number
of available primitive coordinates.  If the final count is still not \(3N-6\)
or \(3N-5\), the run stops.  A downstream solver must never repair the basis by
inventing an independent coordinate model.

\section{Coordinate Definition Format}

Coordinate definitions are printed in a Gaussian ReadGIC-compatible form.  Each
line has a name on the left-hand side, an equal sign, and either one primitive
or a normalized linear combination of primitives on the right-hand side:

\begin{lstlisting}[style=merlino]
Str0001 = R(1,2)
Bend0007 = A(2,1,6)
Tors0014 = D(3,4,5,6)
OuPl0003 = U(7,2,1,3)
GIC0021 =  0.70710678 A(2,1,6) -0.70710678 A(3,4,5)
\end{lstlisting}

\noindent
Atom indexes are one-based and always refer to the input atom order after any
explicit canonical ring ordering has been recorded in the report.  The
primitive symbols are:

\begin{longtable}{@{}p{0.17\linewidth}p{0.73\linewidth}@{}}
\toprule
Primitive & Meaning \\
\midrule
\endhead
\texttt{R(i,j)} & Bond stretch between atoms \(i\) and \(j\). \\
\texttt{A(i,j,k)} & Valence angle \(i-j-k\), with \(j\) as the central atom. \\
\texttt{L(i,j,k)} & Linear-bend component for a nearly linear \(i-j-k\)
arrangement. \\
\texttt{D(i,j,k,l)} & Dihedral angle around the central bond \(j-k\). \\
\texttt{U(i,j,k,l)} & Out-of-plane coordinate with a single central atom.  The
Fortran and Python parsers enforce the same central-atom convention. \\
\texttt{ImpD(i,j,k,l)} & Improper-dihedral compatibility form used only with
\texttt{G16} or \texttt{ImpDih}. \\
\bottomrule
\end{longtable}

\noindent
The name prefix indicates the coordinate family and is part of the stable
ordering contract:

\begin{longtable}{@{}p{0.18\linewidth}p{0.72\linewidth}@{}}
\toprule
Prefix & Family \\
\midrule
\endhead
\texttt{Str} & primitive stretch; \\
\texttt{Bend} & exocyclic or final valence-angle coordinate; \\
\texttt{LinB} & linear-bend coordinate; \\
\texttt{Tors} & exocyclic torsion; \\
\texttt{BtFl} & butterfly torsion around a bond shared by two rings; \\
\texttt{CyGNA} & endocyclic valence-angle combination; \\
\texttt{QPck} & active ring puckering amplitude-like coordinate built from
endocyclic torsions; \\
\texttt{PhiP} & active ring puckering phase-like coordinate built from
endocyclic torsions; \\
\texttt{RPck} & inactive radial puckering component used internally or
reported for diagnostics, not part of the active Gaussian block; \\
\texttt{OuPl} & out-of-plane coordinate; \\
\texttt{ImpD} & improper-dihedral compatibility coordinate. \\
\bottomrule
\end{longtable}

\subsection{Ring Coordinate Format}

For a ring with ordered atoms \(a_1,a_2,\ldots,a_n\), \program{} generates two
ring families.  Endocyclic valence-angle coordinates are combinations of
primitive angles
\[
  A(a_{i-1},a_i,a_{i+1}),
\]
where indexes are cyclic.  Endocyclic torsional coordinates are combinations
of primitive dihedrals
\[
  D(a_{i-1},a_i,a_{i+1},a_{i+2}).
\]
The primitive torsions are then transformed to active puckering coordinates
\texttt{QPck} and \texttt{PhiP}.  These are preferred in Gaussian and
downstream reports because they describe the independent ring deformation more
compactly than the raw endocyclic torsions.  The inactive \texttt{RPck}
coordinate corresponds to the redundant radial component and is not written as
an active final coordinate.

For a five-membered ring, the report therefore separates the coordinate
construction from the active coordinates:

\begin{lstlisting}[style=merlino]
Endocyclic Dihedral Angles
  5-Membered Ring (  13  12  16   1   2):  2 Dihedral angles

Final ring torsion coordinates
  QPck0001 = ...
  PhiP0001 = ...
\end{lstlisting}

\noindent
For six-membered and larger rings, the same rule is used: the primitive basis
is endocyclic torsional, but the final active coordinates are the independent
\(q\)- and \(\phi\)-like ring deformation coordinates selected by the rank
and symmetry checks.  Fused rings add \texttt{BtFl} butterfly coordinates around
shared bonds.  Butterfly coordinates remain in the torsional family for
pruning because their B rows describe ring-puckering motion.

\section{Coordinate Families}

\program{} builds primitive coordinates from the covalent graph.  Hydrogen bonds
may be detected and reported by higher-level topology utilities, but the
standard GIC graph remains covalent: hydrogen bonds are correction targets, not
ring-closing covalent edges.

\begin{longtable}{@{}p{0.23\linewidth}p{0.67\linewidth}@{}}
\toprule
Family & Policy \\
\midrule
\endhead
Stretch & All bonded pairs are kept as primitive \texttt{R(i,j)} coordinates.
Stretches are not reduced by the final pruning step. \\
Exocyclic bend & Valence angles outside rings are generated around each center
and reduced locally by rank. \\
Linear bend & Linear-angle candidates are generated only when required by local
geometry and are capped per center to avoid unphysical proliferation. \\
Exocyclic torsion & The default torsion generator is the one-dihedral
\texttt{ONEDIH} path.  The keyword \texttt{NOONEDIH} forces the older
alternative for diagnostics. \\
Ring bend & Endocyclic valence-angle combinations are generated ring by ring,
with canonical ring numbering. \\
Ring torsion & Endocyclic torsions are converted to ring deformation
coordinates.  The Gaussian block uses the active \texttt{QPck} and
\texttt{PhiP} coordinates; inactive radial \texttt{RPck} components are not
part of the final active basis. \\
Butterfly & Fused-ring butterfly coordinates are generated around shared bonds
and kept as torsional ring-family coordinates. \\
Out-of-plane & Default non-planarity primitives are \texttt{U(...)} out-of-plane
angles.  They are not generated when all four involved atoms are cyclic.  Cyclic
centers with exocyclic substituents remain eligible. \\
Improper dihedral & With \texttt{G16} or \texttt{ImpDih}, improper dihedrals are
written instead of out-of-plane angles.  This is a compatibility mode and must
match the Python parser exactly. \\
\bottomrule
\end{longtable}

\section{Rings and Polycyclic Systems}

Ring atom ordering is canonicalized with the same Prelog-first convention used
by the Python implementation.  The \provout{} report lists atoms in one cycle,
atoms linking multiple cycles, rings found versus expected rings, endocyclic
valence-angle coordinates, endocyclic torsional coordinates, butterfly
coordinates and out-of-plane candidates.

For planar polycyclic aromatic hydrocarbons, out-of-plane candidates whose four
atoms are all cyclic are suppressed.  Exocyclic out-of-plane coordinates remain
available and are either kept or removed by the final rank test.  Regression
targets should include naphthalene, anthracene, pyrene and coronene, with
explicit checks that fused-ring butterfly coordinates are present when
topologically required.

\section{Redundancy Removal}

The final pruning step is type-local.  \program{} builds B rows for candidate
GICs and applies a modified Gram--Schmidt rank test separately to homogeneous
coordinate families:

\begin{itemize}
  \item exocyclic bends;
  \item linear bends;
  \item exocyclic torsions;
  \item butterfly and ring torsional coordinates;
  \item ring valence-angle coordinates;
  \item out-of-plane coordinates.
\end{itemize}

\noindent
Stretch rows are always retained as primitive coordinates.  This policy avoids
the chemically undesirable behavior of a global SVD in which, for example, a
stretch can eliminate a torsion simply because it appears earlier in the global
rank basis.  The rank test is still numerical and B-matrix based, but the
admissible competition is restricted to coordinates with the same physical
role.

The target rank is \(3N-6\) for nonlinear molecules and \(3N-5\) for linear
molecules.  If the final non-redundant count does not match the target,
\program{} stops with an error.  If the pre-pruning count differs from the
target, coordinate definitions are not printed as final coordinates until after
pruning has produced the usable basis.

\section{Symmetry}

\program{} uses symmetry in three different places, and these roles should not
be confused.

\subsection{Point-group perception}

The first role is molecular point-group perception.  \program{} links the
shared Fortran symmetry engine \texttt{symm.f}.  Cartesian coordinates are
centered at the center of mass, oriented on principal inertia axes and assigned
a point group.  The authoritative report is the point group computed by
\texttt{symm.f}; an optional point-group token on the second \xyzin{} line is
legacy metadata only.

The symmetry report records the detected point group, the quality of the
assignment and the maximum atom-matching deviation.  A high-symmetry assignment
is used only when atom permutations and Cartesian operations are consistent
within the configured tolerances.  If the geometry is only approximately
symmetric, the report must make that explicit, because the symmetry-adapted
coordinate basis is only as meaningful as the atom mapping used to build it.

\subsection{Symmetry labels for final GICs}

The second role is assigning irreducible representations to the final
post-pruning GIC basis.  This is done after the non-redundant coordinate set has
been selected.  The order matters: symmetry is not allowed to rescue a
redundant coordinate set, and pruning is not allowed to destroy the final
irreducible-representation counts.  The target representation is
\[
  \Gamma_\mathrm{vib} = \Gamma_{3N}-\Gamma_\mathrm{trans}-\Gamma_\mathrm{rot}.
\]
For a nonlinear molecule, the sum of the target irrep counts is therefore
\(3N-6\).

With \texttt{GICSYM}, the deterministic Python post-check projects the final
GIC rows onto symmetry-adapted combinations, assigns irreducible
representations, writes the symmetry-adapted GIC block and records the result
in:

\begin{itemize}
  \item \texttt{gauin.symm};
  \item \texttt{gicsym};
  \item \texttt{gic\_symmetry\_diagnostics.json}.
\end{itemize}

\noindent
The promoted \gauin{} file contains the symmetrized Gaussian block.  Totally
symmetric coordinates are written first, followed by a blank separator and then
the other irreducible representations.  Downstream programs must not
re-symmetrize; they must consume the frozen schema and its labels.

\subsection{Allowed symmetry blocks}

The post-check is deliberately constrained.  A symmetry-adapted coordinate must
be expressible from a physically homogeneous source block.  The allowed blocks
are primitive stretches, ordinary bends, linear bends, torsions, out-of-plane
coordinates, ring bends, ring torsions and local out-of-plane/improper blocks
around the same center.  There is no global fallback that mixes every primitive
in the molecule.  This preserves the meaning of the final coordinates and makes
the result comparable between Fortran and Python.

Within a homogeneous block, equivalent primitives are collected by the symmetry
operations.  The totally symmetric combination is the normalized sum over the
orbit.  Non-totally-symmetric combinations are orthonormal differences or the
corresponding projection-operator rows.  The sign is canonicalized so repeated
runs produce the same labels and coefficients.  The post-check then verifies
both irrep counts and coordinate-family counts.  For example, a missing
torsional \(B_{2g}\) coordinate cannot be replaced by an extra stretch of the
same irrep.

\subsection{Fortran local SALCs versus production symmetry}

The keyword \texttt{SYMMALL} activates a limited Fortran-side local SALC pass
for selected non-stretch one-term coordinates.  This is a construction aid, not
the owner of production symmetry.  Production symmetry adaptation is the
\texttt{GICSYM} post-check, because it sees the final pruned basis, can compare
against the theoretical vibrational representation and writes the frozen
schema.  Stretches remain primitive in the raw \gauin{} and are symmetry-adapted
only in the post-check when \texttt{GICSYM} is requested.

\subsection{Symmetry-adapted Cartesians}

With \texttt{SYCART}, \program{} writes symmetry-adapted Cartesian coordinates
after removal of rotations and translations.  This provides a common service
for SEfit, MORPHEUS and future optimizers without embedding symmetry logic in
those modules.  The Cartesian basis follows the same point-group assignment as
the GIC basis.  It is useful when a downstream module wants a translationally
and rotationally clean totally symmetric displacement space without using GICs.

\subsection{Failure policy for symmetry}

Symmetry adaptation fails if the projected basis cannot reproduce the target
vibrational irrep counts, if coordinate-family counts change, if atom
permutations are inconsistent, or if two equivalent source blocks would produce
ambiguous labels.  These failures are intentional.  A wrong or ambiguous
symmetry label would contaminate downstream refinement, Gaussian optimization
or force-field projection more seriously than a clean stop.

\section{Frozen Schema and Public API}

The production public API is the \texttt{GICForge} service in
\texttt{merlino\_gic}.  A fresh definition run writes a
\texttt{merlino.gic.definition.v1} schema containing:

\begin{itemize}
  \item primitive coordinate expressions;
  \item GIC coefficient matrix;
  \item coordinate labels and names;
  \item coordinate families;
  \item irreducible representations;
  \item point group and symmetrization flag;
  \item Gaussian-readable GIC block;
  \item provenance hashes for input, backend and generated files.
\end{itemize}

\noindent
The schema is frozen during an ordinary refinement or optimization.  Later
iterations evaluate values and B rows from the schema without changing
topology, ring detection, redundancy removal or symmetry assignment.  A restart
is the explicit boundary at which the schema may be regenerated.

\section{Output Files}

\begin{longtable}{@{}p{0.33\linewidth}p{0.57\linewidth}@{}}
\toprule
File & Meaning \\
\midrule
\endhead
\provout{} & Human-readable report: point group, topology summary, candidate
families, pruning diagnostics, final counts and final coordinate definitions. \\
\gauin{} & Gaussian ReadGIC input block for the final post-pruning coordinate
basis.  With \texttt{GICSYM}, this is the promoted symmetrized block. \\
\texttt{gauin.raw} & Preserved unsymmetrized Gaussian block when \texttt{GICSYM}
promotes \texttt{gauin.symm}. \\
\texttt{gauin.symm} & Symmetry-adapted Gaussian block before promotion. \\
\gicsym{} & Names, irreps and source-block information for the symmetrized
coordinates. \\
\texttt{bmat.out} & Machine-readable final B matrix when \texttt{BMAT} is
active. \\
\texttt{bondout} & Detected covalent topology. \\
\texttt{ringin} & Canonical ring data used by the ring-coordinate generator. \\
\texttt{sycart.xyz} & Symmetry-adapted Cartesian coordinates when
\texttt{SYCART} is active. \\
\texttt{merlino.gic.definition.}\newline\texttt{v1.json} & Frozen reusable GIC schema written by
\texttt{gic-define}. \\
\bottomrule
\end{longtable}

\section{Keyword Summary}

\begin{longtable}{@{}p{0.18\linewidth}p{0.72\linewidth}@{}}
\toprule
Keyword & Effect \\
\midrule
\endhead
\texttt{GNIC} & Build generalized natural/internal coordinates. \\
\texttt{BMAT} & Write the machine-readable final B matrix. \\
\texttt{GICSYM} & Produce symmetry-adapted GICs and promote them to the Gaussian
input block. \\
\texttt{SYCART} & Produce symmetry-adapted Cartesian coordinates without
rotations and translations. \\
\texttt{GDV} & Use out-of-plane \texttt{U(...)} coordinates; this is also the
default behavior. \\
\texttt{G16}, \texttt{ImpDih} & Use Gaussian-style improper dihedrals instead
of out-of-plane coordinates. \\
\texttt{ONEDIH} & One-dihedral torsion generation.  This is the default. \\
\texttt{NOONEDIH} & Force the legacy alternative torsion path for diagnostics. \\
\texttt{SYMMALL} & Apply the limited type-local Fortran SALC generator before
final pruning.  Production symmetry adaptation is still owned by
\texttt{GICSYM}. \\
\bottomrule
\end{longtable}

\section{Python/Fortran Identity Contract}

The Python layer is allowed to orchestrate, serialize, post-check symmetry and
evaluate the frozen B matrix.  It is not allowed to replace a valid Fortran GIC
basis silently.  The contract checks:

\begin{itemize}
  \item final coordinate count and family counts;
  \item ordered primitive signatures;
  \item point group and irreducible-representation labels;
  \item number of totally symmetric coordinates;
  \item Python analytic B matrix against Fortran \texttt{bmat.out};
  \item byte-stable symmetry diagnostics for repeated runs.
\end{itemize}

\noindent
Run the contract with:

\begin{lstlisting}[style=merlino]
python -m merlino gic-contract --geometry examples/pyrene.xyz \
  --workdir work/gic_contract_pyrene
\end{lstlisting}

\noindent
The Python port in \texttt{merlino\_gic.gicforge\_python} is a faithful porting
target, not an independent production generator.  It must match Fortran on
ordering, primitive orientation, coordinate families, type-local pruning,
symmetry labels and B rows before it can replace a backend path.

\section{Recommended Regression Set}

The minimum regression set should cover:

\begin{itemize}
  \item acyclic molecules with ordinary stretches, bends and torsions;
  \item linear-angle cases with capped linear bends;
  \item square-planar or high-coordination centers;
  \item SF\(_6\), which should not produce spurious linear bends;
  \item saccharine and other heteroatom rings;
  \item naphthalene, anthracene, pyrene and coronene for fused PAHs;
  \item norbornene and norbornadiene for bridged non-planar rings;
  \item \texttt{GICSYM} variants for nontrivial point groups;
  \item \texttt{SYCART} variants for symmetry-adapted Cartesian output.
\end{itemize}

\noindent
Every test should verify the target rank, coordinate-family counts, presence of
required butterfly coordinates, absence of all-cyclic out-of-plane primitives,
and equality of Python/Fortran B matrices where both paths are active.

\section{Failure Policy}

\program{} must fail early when the generated basis is not scientifically
usable.  Fatal conditions include:

\begin{itemize}
  \item inconsistent symbol/atomic-number input;
  \item topology changes during a frozen-schema workflow;
  \item final GIC count different from \(3N-6\) or \(3N-5\);
  \item symmetry post-check unable to reach theoretical irrep counts;
  \item coordinate-family counts changed by post-check symmetrization;
  \item Python/Fortran B-matrix mismatch beyond numerical roundoff.
\end{itemize}

\noindent
These failures are intentional.  A downstream fit or optimizer must not proceed
with an ambiguous coordinate definition.

\section{Relation to MORPHEUS and SEfit}

MORPHEUS and SEfit call \program{} as a coordinate-service library.  A fresh fit
generates one frozen coordinate definition at the starting geometry.  Ordinary
least-squares iterations then reuse that definition and only update coordinate
values, B rows and projected parameter directions.  This prevents the fitted
model from changing during optimization and makes the final refinement
auditable.

For a GIC active model, the fit uses the totally symmetric block assigned by
\program{}.  For a symmetry-adapted Cartesian active model, the fit uses
\texttt{SYCART}.  In both cases, topology and symmetry belong to
\program{}, while weights, constraints, predicates, parameter classes and
least-squares diagnostics belong to the refinement layer.

\end{document}
